\section*{第 \ref{ch:b3:cell-cycle-apoptosis} 章——细胞周期调控与程序性细胞死亡}

\begin{solution}{exo:b3:cell-cycle-apoptosis:1}
G1、S、G2、M。G1 期推进与\omterm{def:b3:cell-cycle-apoptosis:checkpoints}{限制点}：周期蛋白 D--Cdk4/6；
G1/S：周期蛋白 E--Cdk2；S：周期蛋白 A--Cdk2；G2/M：周期蛋白 B--Cdk1。
\omterm{def:b3:cell-cycle-apoptosis:cdk}{后期促进复合物}（APC/C）降解周期蛋白 B 和\omterm{prop:b3:cell-cycle-apoptosis:anaphase}{分离酶抑制蛋白}，
从而结束有丝分裂。
\end{solution}

\begin{solution}{exo:b3:cell-cycle-apoptosis:2}
哈特韦尔：出芽酵母的 \emph{cdc} 突变体，停滞于特定
阶段，以及“Start”的概念。纳斯：裂殖酵母的 \emph{cdc2}，
作为调控有丝分裂时机的激酶，其人类同源物（Cdk1）能够挽救
酵母缺陷——体现保守性。亨特：海胆卵中的周期蛋白，
一种在整个周期中持续合成、并在每次分裂时被降解的蛋白质。
\end{solution}

\begin{solution}{exo:b3:cell-cycle-apoptosis:3}
\omterm{def:b3:cell-cycle-apoptosis:checkpoints}{限制点}（G1 晚期）：生长因子、营养物质、细胞大小；作用于
周期蛋白 D--Cdk4/6 和 Rb。G2/M：DNA 完整且已复制；\omterm{def:b3:dna-repair:ddr}{ATM}/ATR
抑制 Cdc25，使 Cdk1 保持磷酸化而失活。\omterm{def:b3:cell-cycle-apoptosis:checkpoints}{纺锤体
组装检查点}：每个动粒均已附着；未附着的动粒会形成
抑制 APC/C 的 Mad2--Cdc20 复合体。
\end{solution}

\begin{solution}{exo:b3:cell-cycle-apoptosis:4}
\omterm{def:b3:cell-cycle-apoptosis:apoptosis}{细胞凋亡}：细胞皱缩，\omterm{def:b3:chromatin-epigenetics:nucleosome}{染色质}凝聚，DNA 被切割成梯状条带，
膜出现泡状突起但保持封闭，\omterm{def:b3:cell-cycle-apoptosis:apoptosis}{磷脂酰丝氨酸}外露，尸体
被吞噬，不引发炎症。\omterm{def:b3:cell-cycle-apoptosis:apoptosis}{坏死}：细胞肿胀，膜破裂，
内容物外漏，DNA 随机降解，引发炎症。执行型：
\omterm{def:b3:cell-cycle-apoptosis:apoptosis}{胱天蛋白酶}，是一类在天冬氨酸之后切割的半胱氨酸蛋白酶。
\end{solution}

\begin{solution}{exo:b3:cell-cycle-apoptosis:5}
上升段由 $5$ 升至 $30$，速率为每分钟 $1$，耗时 \qty{25}{min}。降解段由
$30$ 降至 $5$，半衰期为 \qty{2}{min}：$\log_{2}(30/5)\times 2 =
\qty{5.2}{min}$。周期约为 \qty{30}{min}，有丝分裂占其 \qty{17}{\%}。
\end{solution}

\begin{solution}{exo:b3:cell-cycle-apoptosis:6}
(a) $C$ 永远不会降到 $C_{1}$ 以下，细胞被锁定于高值分支、停留在有丝分裂期。
(b) 无抑制性磷酸化：阈值 $C_{2}$ 降低，
有丝分裂在细胞较小时即开始（即“wee”表型）。
(c) Cdk1 保持磷酸化：高值分支无法到达，细胞停滞于
G2 期，细胞拉长。(d) 与 (b) 相同，还有 G2/M \omterm{def:b3:cell-cycle-apoptosis:checkpoints}{检查点}
经由 Wee1/Cdc25 起作用，已无法再阻止该细胞。
\end{solution}

\begin{solution}{exo:b3:cell-cycle-apoptosis:7}
置于 S 期细胞质中的 G1 期细胞核：其复制起点已获\omterm{def:b3:cell-cycle-apoptosis:checkpoints}{复制许可}，且
细胞质提供有活性的周期蛋白 E/A--Cdk2，因此立即进入 S 期。G2
置于 S 期细胞质中的 G2 期细胞核：其复制起点已启动，且未重新获得许可
（\omterm{def:b3:cell-cycle-apoptosis:checkpoints}{复制许可}因子被破坏或被 CDK 抑制），因此
不能再次复制；它等待伙伴细胞也达到 G2 期，两者
共同进入有丝分裂。
\end{solution}

\begin{solution}{exo:b3:cell-cycle-apoptosis:8}
未附着的动粒是一种催化剂：它不断地把 Mad2 转变为
其活性构象，由此产生一种可扩散的
Cdc20 抑制剂，它扩散至整个细胞，使 APC/C 处处保持关闭
——一个动粒即可产生足够的抑制剂。无 Mad2 时，APC/C 一经 Cdk1 开启便被激活，
不论染色体是否已附着；
后期在染色体未附着时即已开始，子细胞
成为非整倍体，生物体（或细胞系）因染色体丢失而死亡。
\end{solution}

\begin{solution}{exo:b3:cell-cycle-apoptosis:9}
缺失 \emph{ced-9}：本应存活的细胞死亡——胚胎致死，
原因是死亡过多。缺失 \emph{ced-3}：这 $131$ 次死亡均不发生，
细胞存活并分化，线虫能够存活。两者同时缺失：无死亡，
——即 \emph{ced-3} 表型。由于去除执行者会抵消去除守护者的效应，
CED-9 作用于上游，通过维持
CED-4/CED-3 处于非活性状态而发挥作用；失去 CED-9 时它们即被激活，除非
它们本身也缺失。
\end{solution}

\begin{solution}{exo:b3:cell-cycle-apoptosis:10}
令 $A = aC$，$\mathrm{d}C/\mathrm{d}t = k_{s} - k_{d}aC^{2}$，这是一个
单变量方程，其稳态为 $C^{*} = \sqrt{k_{s}/(k_{d}a)}$，
该处斜率为 $-2k_{d}aC^{*} < 0$：任何偏离都会单调衰减，
单一一阶变量不能振荡。
S 形曲线给出两条稳定分支，二者之间隔着禁区，
因此状态必然跳变而无法稳定下来；若反馈中存在显式延迟，
则会以另一种途径产生振荡——制动作用
来得太晚、出现过冲等——许多激素
节律正是如此。
\end{solution}

\begin{solution}{exo:b3:cell-cycle-apoptosis:11}
每个到达的感受器都被一个守护者捕获，因此在
所有 $N$ 个守护者被占据之前什么都不会发生，此时 $t \approx N/r$；随后的
感受器不再受约束，激活 Bax/Bak，其成孔随数量急剧增加，
线粒体在数分钟内开放。更强的
刺激只能通过 $r$ 缩短时间；超过阈值后
结局相同。Venetoclax 占据守护者，使
有效 $N$ 减小：时间缩短，而已接近 $N$ 的致敏细胞——
即刻死亡。
\end{solution}

\begin{solution}{exo:b3:cell-cycle-apoptosis:12}
过量的 Bcl-2 阻断内源性途径：本应在生长信号消失后死亡的 B~细胞
得以存活，细胞群体因死亡失效而增长，
增长速率缓慢，取决于此类细胞的产生速率——惰性的，
直至第二个\omterm{def:b3:dna-repair:lesions}{损伤}带来增殖。Rb 缺失使
\omterm{def:b3:cell-cycle-apoptosis:checkpoints}{限制点}失效：视网膜前体细胞可以无需生长因子而
通过\omterm{def:b3:cell-cycle-apoptosis:cdk}{细胞周期}，并以细胞分裂机器所允许的最快速度分裂——
侵袭性的。两种肿瘤分别对应两个程序各自单独失效：
即死亡控制与分裂控制，二者中任一缺失都足以形成肿瘤，
后者进展更快。
\end{solution}

\begin{solution}{pb:b3:cell-cycle-apoptosis:1}
\textbf{1.} $10^{7}\times 3500/5 = 7\times 10^{9}$ 个细胞，每天约
每秒约 \num{80000} 个。
\textbf{2.} 每个隐窝每天 $3500/5 = 700$；$22\times 2^{5} = 704$。
\textbf{3.} 约 $80\times 365 \approx \num{29000}$ 次分裂；$\times 0.6
\approx \num{17500}$ 个突变。
\textbf{4.} 过渡细胞数日内脱落，
连同其中的突变一并移出；干细胞终生存在，故其
每一个突变都被保留，且每次分裂又增添新突变——缓慢循环
可将其减至最少。
\textbf{5.} 绒毛约在其正常过渡的 $5$ 天内排空；
重建需要一天的恢复期，五次分裂
每次 \qty{12}{h}，并沿绒毛迁移——共四至五天，
此间屏障裸露：肠道快速过渡的分裂使其
成为辐射的早期受害者。
\textbf{6.} 密封的凋亡细胞残体不向充满细菌的管腔释放内容物，
不引发炎症；\omterm{def:b3:cell-cycle-apoptosis:apoptosis}{坏死}则会溢出
酶和信号，每脱落一个细胞都会引起黏膜炎症。
\textbf{7.} $(40 - 8)/2 = \qty{16}{h}$。
\textbf{8.} $\log_{2}(40/8)\times \qty{10}{min} = 2.3\times 10 =
\qty{23}{min}$。
\textbf{9.} 周期约 \qty{16.4}{h}；Cdk1 有活性的
时间占 \qty{2.3}{\%}。
\textbf{10.} 不仅取决于 $k_{s}$，更取决于在阈值处停留的时间：
体细胞在\omterm{def:b3:cell-cycle-apoptosis:checkpoints}{限制点}等待生长因子，并在
G2/M 等待其 DNA 完成，干细胞停留的时间长于过渡细胞。
蛙卵没有\omterm{def:b3:cell-cycle-apoptosis:checkpoints}{检查点}，细胞质中储备了一切所需物质，
因此仅由裸振子决定其周期。
\textbf{11.} 位于低分支上，周期蛋白水平高于正常的
$C_{2}$ 值：\omterm{def:b3:cell-cycle-apoptosis:checkpoints}{检查点}通过抑制 Cdc25 使 S 形曲线右移，
因而跃变不会发生。修复后 Cdc25 被释放，
阈值降至已积累的周期蛋白水平之下，
细胞立即进入有丝分裂——这正是\omterm{def:b3:cell-cycle-apoptosis:checkpoints}{检查点}解除后的细胞
同步进入有丝分裂的原因。
\textbf{12.} 周期蛋白 B 与\omterm{prop:b3:cell-cycle-apoptosis:anaphase}{分离酶抑制蛋白}均未降解：细胞停留在
中期，\omterm{prop:b3:cell-cycle-apoptosis:anaphase}{黏连蛋白}完好，Cdk1 保持活性，无限期维持；细胞在此死亡，
或最终带着未分离的基因组逸出。
\textbf{13.} 每天 $10^{11}$ 次分裂。
\textbf{14.} 有\omterm{def:b3:cell-cycle-apoptosis:checkpoints}{检查点}时每天 $10^{11}/5\times 10^{4} = 2\times 10^{6}$
个非整倍体子细胞；无\omterm{def:b3:cell-cycle-apoptosis:checkpoints}{检查点}时为 $2\times 10^{9}$。
\textbf{15.} 每天 $2\times 10^{4}$；历经 $80$ 年共 $6\times 10^{8}$。
\textbf{16.} 每天 $10^{9}/50 = 2\times 10^{7}$ 次错误分离：每
天肿瘤都会尝试两千万种新核型，选择保留
生长更快或耐药的那些。
\textbf{17.} 若每次分裂都发生错误分离，则
任何胚胎细胞谱系都无法保持整倍体基因组，胚胎死亡。部分丢失则
偶发的非整倍性见于存活下来的细胞（尤其是在缺失
\omterm{def:b3:dna-repair:ddr}{p53} 时），这是一种为肿瘤提供变异、
而不致杀死生物体的染色体不稳定性。
\textbf{18.} 减数分裂错误将多余的染色体送入配子，
因此受精卵及由其衍生的每个细胞均为三体；有丝分裂
错误发生于单个体细胞，仅由其克隆继承。
\textbf{19.} $10^{11}\times \qty{1}{ng} = \qty{100}{g}$ 每天，
每年 \qty{36}{kg}，约相当于机体质量的一半，每年如此。
\textbf{20.} 每天 \qty{20}{g} 蛋白质，约占膳食摄入量的三分之一，
约占机体蛋白质总周转量的百分之几。
\textbf{21.} $10^{11}$ 个凋亡残体，每个巨噬细胞每天处理 $24$ 个：$4\times
10^{9}$ 个巨噬细胞全时工作，占 \qty{4}{\%} 的 $10^{11}$
个巨噬细胞的时间。
\textbf{22.} 裂解的残体释放蛋白酶、DNA、ATP 及其他
警报信号，引发炎症，并可能诱发
针对核抗原的自身免疫；完整残体则
借其外小叶上的\omterm{def:b3:cell-cycle-apoptosis:apoptosis}{磷脂酰丝氨酸}（并丢失“不要吃我”信号）
标示自身，并以释放的核苷酸吸引吞噬细胞。
\textbf{23.} 内源性途径在线粒体处受阻。
外源性途径仍然有效，胱天蛋白酶-8 可直接激活胱天蛋白酶-3，
但经由 Bid 的放大作用已丧失。细胞
积累是由于不死亡而非分裂更快——即缓慢生长
——直到后续\omterm{def:b3:dna-repair:lesions}{损伤}（常为 \emph{MYC}）带来增殖。
\textbf{24.} \omterm{def:b3:cell-cycle-apoptosis:apoptosis}{细胞凋亡}需要数小时，依赖 ATP，并经过
可被阻断的步骤——\omterm{def:b3:cell-cycle-apoptosis:apoptosis}{胱天蛋白酶}抑制剂、线粒体
阈值——因此半影区神经元若及时处理可被挽救；
核心区神经元则在数分钟内因能量衰竭而死亡，
没有可供中断的程序。
\textbf{25.} 裸振子周期约 \qty{16}{h}；正常机体中每天约有 $2\times
10^{4}$ 个非整倍体分裂细胞；每天通过\omterm{def:b3:cell-cycle-apoptosis:apoptosis}{细胞凋亡}回收
\qty{100}{g} 细胞。
\end{solution}
